% Transcribed from CMM110 Midterm Review Question.docx.
\IfFileExists{Book\space NL/template/kaobook.cls}{
  \IfFileExists{main.tex}{\documentclass[main.tex]{subfiles}}{
    \IfFileExists{../../main.tex}{\documentclass[../../main.tex]{subfiles}}{
      \IfFileExists{../main.tex}{\documentclass[../main.tex]{subfiles}}{
        \documentclass[Book\space NL/main.tex]{subfiles}
      }
    }
  }
}{
  \IfFileExists{../../template/kaobook.cls}{\documentclass[../../main.tex]{subfiles}}{\documentclass{book}}
}
\IfFileExists{Book\space NL/template/kaobook.cls}{}{\IfFileExists{../../template/kaobook.cls}{}{\InputIfFileExists{review-notes-standalone.tex}{}{\InputIfFileExists{Book\space NL/appendices/appendix 4-Review Notes/review-notes-standalone.tex}{}{}}}}
\begin{document}
\ifSubfilesClassLoaded{
  \ifcsname appendixchapterparent\endcsname
  \else
    \standaloneappendixchaptersetup{D}{3}
  \fi
}{}
\IfFileExists{\booknlroot appendices/appendix 4-Review Notes/images/d3-figure-aquaporin-localization-kidney-cells.png}{
  \edef\reviewnotesimagepath{\booknlroot appendices/appendix 4-Review Notes/images/}
}{
  \edef\reviewnotesimagepath{\booknlroot appendices/appendix 4-Review Notes/images/}
}
\chapter{Cell Biology Review Questions}
\labch{appendix-d3-cellular-molecular-biology-review-questions}

\margintoc
\section{Topics Listed Without Review Questions}
The supplied review document names these topics but provides no accompanying questions or study notes:
\begin{itemize}
\item Cell diversity and visualization.
\item Cellular organelles.
\item Protein sorting and intracellular trafficking.
\item Cytoskeleton and cell motility.
\item Cell signaling and communication.
\item Cell growth and cancer.
\item Cell injury and cell death.
\item Stem cells and tissue regeneration.
\item Biomolecules.
\item Enzymes and cofactors.
\item Bioenergetics and cellular respiration.
\end{itemize}
\section{Membrane Structure and Function}
\subsection{Plasma Membrane Structure}
\begin{enumerate}
\item How do phospholipid transporters contribute to the negative charge on the cytosolic face of the plasma membrane?
\item You are studying a protein involved in signaling in response to hormone treatment in cells. Subcellular fractionation experiments indicate that your protein is associated with the plasma membrane. How can you determine whether it’s a peripheral or integral membrane protein?
\item You raise antibodies against two peptides of a plasma membrane-associated protein, one from its N-terminus and the other from its C-terminus. When you use the antibodies in immunofluorescence, you find that the antibody to the N-terminus stains living cells, but the antibody to the C-terminus only stains fixed cells. Diagram the structure of the protein in the plasma membrane.
\item How did the cell fusion experiments by Frye and Edidin support the fluid mosaic model of membrane structure? What results would they have obtained if they had incubated fused cells at 2°C?
\item How would treatment of cells with a drug that disrupts actin filaments affect the mobility of plasma membrane proteins?
\item What are lipid rafts, and what cellular processes are they involved in?
\end{enumerate}
\begin{center}\includegraphics[width=0.82\textwidth,height=0.68\textheight,keepaspectratio]{\reviewnotesimagepath d3-figure-aquaporin-localization-kidney-cells.png}\par\captionof{figure}{Immunofluorescence micrographs comparing aquaporin localization in untreated and ADH-treated kidney cells.}\labfig{appendix-figure-d3-aquaporin-localization-adh-immunofluorescence}\end{center}
\begin{center}\includegraphics[width=0.82\textwidth,height=0.68\textheight,keepaspectratio]{\reviewnotesimagepath d3-figure-cftr-localization-wild-type-mutant.png}\par\captionof{figure}{Electron micrographs comparing mock-infected cells, wild-type CFTR localization, and the delta-F508 CFTR mutant.}\labfig{appendix-figure-d3-cftr-wild-type-mutant-localization}\end{center}
\subsection{Membrane Transport}
\begin{enumerate}
\item The concentration of K+ is about 20 times higher inside squid axons than in extracellular fluids, generating an equilibrium membrane potential of –75 mV. What would be the expected equilibrium membrane potential if the K+ concentration were only ten times higher inside than outside the cell?\\ \emph{(Given: R = 1.98 × 10⁻³ kcal/mol/deg, T = 298° K, ln(x) = 2.3 log₁₀(x), F = 23 kcal/V/mol.)}
\item Curare binds to nicotinic acetylcholine receptors and prevents them from opening. How would this affect muscle contraction?
\item How can K+ channels selectively transport K+ ions but not Na+ ions, given that Na+ ions are smaller?
\item Why is the Na⁺-K⁺ pump not necessary to prevent osmotic lysis in plant cells?
\item How would an inhibitor of ABC transporters affect the drug sensitivity of cancer cells?
\item How can glucose be transported against its concentration gradient without direct ATP expenditure in intestinal epithelial cells?
\end{enumerate}
\subsection{Endocytosis}
\begin{enumerate}
\item How would an inhibitor of actin polymerization affect phagocytosis? Would it also affect clathrin-mediated endocytosis?
\item How are clathrin-coated vesicles targeted to early endosomes?
\item What experiment could you perform to determine whether cells have a limited number of specific binding sites for a macromolecule?
\item How would an inhibitor of the H⁺ pump affect receptor recycling in endocytosis?
\end{enumerate}
\section{Cell Biology: Multiple-Choice Review}
\begin{enumerate}
\item \textbf{All of the following are true about eukaryotic cells EXCEPT:}
\begin{itemize}
\item They have linear chromosomes enclosed in a nuclear membrane.
\item They contain mitochondria for energy production.
\item They have a higher surface-to-volume ratio compared to prokaryotes.
\item Their ribosomes are larger than those in prokaryotes.
\item Their membranes are composed of phospholipids linked by ester bonds.
\end{itemize}
\item \textbf{Which of the following best explains how super-resolution microscopy improves upon confocal microscopy?}
\begin{itemize}
\item By using advanced objective lenses with higher numerical apertures.
\item By employing photoactivation to selectively illuminate specific molecules.
\item By eliminating fluorescence from outside the focal plane with a pinhole aperture.
\item By extending the depth of focus through structured illumination.
\item By using multiple wavelengths to create higher contrast images.
\end{itemize}
\item \textbf{A 5-year-old boy presents with skeletal abnormalities and recurrent infections. His lysosomal enzymes are found in elevated levels in the blood. Which of the following cellular processes is most likely disrupted?}
\begin{itemize}
\item Protein degradation within lysosomes.
\item Secretion of proteins through the Golgi apparatus.
\item Transport of enzymes to the mitochondria.
\item Endocytosis of extracellular materials.
\item Assembly of ribosomal subunits in the nucleolus.
\end{itemize}
\item \textbf{Which of the following best describes the role of the proteasome in cellular function?}
\begin{itemize}
\item Lipid biogenesis
\item Degradation of misfolded or damaged proteins
\item Assembly of polypeptide chains
\item Phosphorylation of enzymes
\item Storage of calcium ions
\end{itemize}
\item \textbf{In a phosphoglyceride, which component is responsible for the molecule's hydrophilic properties?}
\begin{itemize}
\item Glycerol backbone
\item Fatty acid tails
\item Phosphate group
\item Choline group
\item Hydrocarbon chains
\end{itemize}
\item \textbf{How do peripheral membrane proteins typically interact with the lipid bilayer?}
\begin{itemize}
\item By spanning the bilayer with hydrophobic domains
\item By attaching to the membrane through electrostatic interactions with lipid head groups
\item By forming covalent bonds with fatty acids
\item By embedding in the lipid core through alpha helices
\item By transiently interacting with transmembrane proteins
\end{itemize}
\item \textbf{Which of the following modifications occurs within the Golgi apparatus and is essential for proper protein trafficking?}
\begin{itemize}
\item Formation of disulfide bonds
\item N-linked glycosylation of proteins
\item Phosphorylation of mannose residues
\item Ubiquitination of misfolded proteins
\item Cleavage of signal peptides
\end{itemize}
\item \textbf{How many of the following statements about the endoplasmic reticulum (ER) are true?} i. It is involved in the synthesis of membrane lipids. ii. The rough ER is involved in protein folding and quality control. iii. Lysosomal enzymes are synthesized in the rough ER. iv. Steroid hormone synthesis occurs in the smooth ER.
\begin{itemize}
\item 4
\item 2
\item 3
\item 1
\item 0
\end{itemize}
\item \textbf{Which cellular process is primarily affected by defects in the function of actin filaments?}
\begin{itemize}
\item Organelle transport
\item Maintenance of cell shape and mechanical support
\item Nuclear envelope breakdown during mitosis
\item Organization of microtubules in the spindle apparatus
\item Mitochondrial division
\end{itemize}
\item \textbf{What is the primary role of formins in actin filament dynamics?}
\end{enumerate}
\begin{itemize}
\item To stabilize existing filaments
\item To promote nucleation and polymerization of actin filaments
\item To sever actin filaments and depolymerize them
\item To initiate the formation of branched actin networks
\item To prevent the binding of ATP to actin monomers
\end{itemize}
\begin{enumerate}
\item \textbf{Which of the following represents a slow cellular response to a growth factor signal?}
\end{enumerate}
\begin{itemize}
\item Activation of ion channels leading to immediate changes in membrane potential.
\item Degradation of proteins via the ubiquitin-proteasome pathway.
\item Activation of transcription factors and changes in gene expression.
\item Rapid polymerization of actin filaments.
\item Opening of calcium channels in the endoplasmic reticulum.
\end{itemize}
\begin{enumerate}
\item \textbf{Which type of receptor is involved in the intracellular signaling of lipid-soluble hormones like cortisol?}
\end{enumerate}
\begin{itemize}
\item G protein-coupled receptor
\item Receptor tyrosine kinase
\item Ligand-gated ion channel
\item Nuclear hormone receptor
\item Toll-like receptor
\end{itemize}
\begin{enumerate}
\item \textbf{What is the likely outcome if meiosis occurs without proper chromosomal segregation?}
\end{enumerate}
\begin{itemize}
\item Cells will undergo apoptosis.
\item Cells will contain abnormal numbers of chromosomes.
\item Cells will contain multiple nuclei.
\item Cells will halt in the G1 phase of the cell cycle.
\item Cells will not undergo cytokinesis.
\end{itemize}
\begin{enumerate}
\item \textbf{How do cancer cells bypass normal cell cycle checkpoints to sustain uncontrolled growth?}
\end{enumerate}
\begin{itemize}
\item By inhibiting enzymes involved in DNA repair.
\item By upregulating apoptosis-related proteins.
\item By activating proto-oncogenes and inactivating tumor suppressors.
\item By maintaining normal levels of mitotic arrest.
\item By downregulating growth factor receptors.
\end{itemize}
\begin{enumerate}
\item \textbf{What does metaplasia refer to?}
\end{enumerate}
\begin{itemize}
\item Accumulation of abnormal proteins in cells.
\item Increase in cell size due to increased protein synthesis.
\item Replacement of one differentiated cell type with another.
\item Formation of fluid-filled vesicles within a tissue.
\item Increase in the number of cells through mitosis.
\end{itemize}
\begin{enumerate}
\item \textbf{Which of the following is NOT typically a characteristic of apoptosis?}
\end{enumerate}
\begin{itemize}
\item DNA fragmentation
\item Formation of apoptotic bodies
\item Cellular swelling and rupture
\item Activation of caspases
\item Plasma membrane blebbing
\end{itemize}
\begin{enumerate}
\item \textbf{Which of the following processes is most crucial for tissue homeostasis and regeneration?}
\end{enumerate}
\begin{itemize}
\item Differentiation of multipotent stem cells into specialized cells.
\item Symmetric division of stem cells to maintain the stem cell pool.
\item Self-renewal of stem cells through mitotic divisions.
\item Conversion of somatic cells into pluripotent stem cells.
\item Fusion of cells to repair damaged tissues.
\end{itemize}
\begin{enumerate}
\item \textbf{How do immune cells, such as T cells, communicate with other cells to coordinate immune responses?}
\end{enumerate}
\begin{itemize}
\item Autocrine signaling
\item Paracrine signaling
\item Endocrine signaling
\item Juxtacrine signaling
\item Neuronal signaling
\end{itemize}
\begin{enumerate}
\item \textbf{Which of the following best describes the structure of proteoglycans in the extracellular matrix?}
\end{enumerate}
\begin{itemize}
\item A core protein attached to repeating units of glycosaminoglycans.
\item A long chain of disaccharides with attached proteins.
\item A membrane-bound receptor linked to glycoproteins.
\item A protein backbone attached to various types of phospholipids.
\item A lipid core surrounded by oligosaccharide chains.
\end{itemize}
\begin{enumerate}
\item \textbf{In the diagram, bonds between various biomolecules are shown. Which bond is primarily involved in stabilizing the secondary structure of proteins?}
\end{enumerate}
\begin{itemize}
\item Peptide bond
\item Hydrogen bond
\item Ionic bond
\item Disulfide bond
\item Covalent bond
\end{itemize}
\begin{enumerate}
\item \textbf{What is the role of a coenzyme in enzymatic reactions?}
\end{enumerate}
\begin{itemize}
\item To provide structural support to the enzyme.
\item To act as a catalyst by lowering the activation energy of the reaction.
\item To donate or accept chemical groups in reactions.
\item To regulate the enzyme's active site through allosteric interactions.
\item To bind substrates and facilitate their conversion to products.
\end{itemize}
\begin{enumerate}
\item \textbf{Which of the following factors is most likely to increase the rate of an enzyme-catalyzed reaction in the Michaelis–Menten model?}
\end{enumerate}
\begin{itemize}
\item Decreasing substrate concentration
\item Increasing product concentration
\item Increasing enzyme concentration
\item Raising the pH beyond the optimal range
\item Adding a non-competitive inhibitor
\end{itemize}
\begin{enumerate}
\item \textbf{Which phosphorylated compound provides the highest energy for substrate-level phosphorylation?}
\end{enumerate}
\begin{itemize}
\item Creatine phosphate
\item ATP
\item GTP
\item 1,3-bisphosphoglycerate
\item Phosphoenolpyruvate
\end{itemize}
\begin{enumerate}
\item \textbf{Which vitamin acts as a coenzyme in redox reactions within the TCA cycle?}
\end{enumerate}
\begin{itemize}
\item Thiamine
\item Biotin
\item Riboflavin
\item Pantothenic acid
\item Niacin
\item \textbf{Which of the following metabolic intermediates is the first committed step in the glycolysis pathway?}
\end{itemize}
\begin{itemize}
\item Glucose 6-phosphate
\item Fructose 1,6-bisphosphate
\item Pyruvate
\item Phosphoenolpyruvate
\item Fructose 6-phosphate
\end{itemize}
\begin{enumerate}
\item \textbf{In oxidative phosphorylation, which of the following molecules serves as the final electron acceptor in the electron transport chain?}
\end{enumerate}
\begin{itemize}
\item NADH
\item Oxygen
\item FADH2
\item Carbon dioxide
\item Water
\end{itemize}
\begin{enumerate}
\item \textbf{Which cellular process is directly affected by a mutation in the gene encoding for RNA polymerase II?}
\end{enumerate}
\begin{itemize}
\item DNA replication
\item Transcription of messenger RNA (mRNA)
\item Splicing of pre-mRNA
\item Translation of mRNA into proteins
\item Protein folding in the endoplasmic reticulum
\end{itemize}
\begin{enumerate}
\item \textbf{What is the primary function of microRNAs (miRNAs) in the regulation of gene expression?}
\end{enumerate}
\begin{itemize}
\item To inhibit the binding of RNA polymerase to DNA
\item To facilitate the translation of specific mRNAs
\item To degrade or prevent the translation of target mRNAs
\item To initiate the synthesis of complementary DNA strands
\item To enhance the binding of ribosomes to mRNA
\end{itemize}
\begin{enumerate}
\item \textbf{Which of the following best describes the role of telomerase in cellular division?}
\end{enumerate}
\begin{itemize}
\item It catalyzes the formation of peptide bonds between amino acids.
\item It prevents the shortening of telomeres by adding repetitive sequences to chromosome ends.
\item It repairs double-strand breaks in DNA during replication.
\item It regulates the disassembly of the mitotic spindle.
\item It initiates the process of apoptosis in damaged cells.
\end{itemize}
\begin{enumerate}
\item \textbf{Which of the following enzymes is responsible for unwinding the DNA double helix during replication?}
\end{enumerate}
\begin{itemize}
\item DNA polymerase
\item Helicase
\item Topoisomerase
\item Ligase
\item Primase
\end{itemize}
\begin{enumerate}
\item \textbf{In the context of cell signaling, which of the following processes is primarily regulated by G protein-coupled receptors (GPCRs)?}
\end{enumerate}
\begin{itemize}
\item Phosphorylation of tyrosine residues
\item Opening of voltage-gated ion channels
\item Activation of adenylyl cyclase and production of cyclic AMP (cAMP)
\item Activation of receptor tyrosine kinases
\item Direct interaction with transcription factors
\end{itemize}
\begin{enumerate}
\item \textbf{Which of the following mechanisms prevents the cell cycle from progressing from G1 to S phase when DNA damage is detected?}
\end{enumerate}
\begin{itemize}
\item Activation of cyclin-dependent kinase (CDK)
\item Degradation of cyclins
\item Inhibition of p53 activity
\item Phosphorylation of retinoblastoma (Rb) protein
\item Activation of DNA repair pathways
\end{itemize}
\begin{enumerate}
\item \textbf{Which type of mutation is most likely to result in a truncated, non-functional protein?}
\end{enumerate}
\begin{itemize}
\item Missense mutation
\item Silent mutation
\item Frameshift mutation
\item Synonymous mutation
\item Inversion mutation
\end{itemize}
\begin{enumerate}
\item \textbf{What is the primary function of the spliceosome during mRNA processing?}
\end{enumerate}
\begin{itemize}
\item To add a poly-A tail to the mRNA
\item To remove introns and join exons together
\item To add a 5' cap to the mRNA
\item To prevent mRNA degradation by nucleases
\item To transport mRNA out of the nucleus
\end{itemize}
\begin{enumerate}
\item \textbf{Which of the following best describes the difference between euchromatin and heterochromatin?}
\end{enumerate}
\begin{itemize}
\item Euchromatin is transcriptionally inactive, while heterochromatin is active.
\item Heterochromatin is found only in the nucleus, while euchromatin is found in the cytoplasm.
\item Euchromatin is loosely packed and associated with active gene expression, while heterochromatin is densely packed and transcriptionally silent.
\item Euchromatin consists of repetitive DNA sequences, while heterochromatin contains protein-coding genes.
\item Euchromatin is associated with DNA repair, while heterochromatin prevents damage.
\end{itemize}
\begin{enumerate}
\item \textbf{Which of the following best describes the function of chaperone proteins in the cell?}
\end{enumerate}
\begin{itemize}
\item To catalyze the folding of proteins into their functional forms
\item To degrade misfolded proteins via the proteasome
\item To bind mRNA and facilitate translation
\item To promote the synthesis of new proteins in the rough ER
\item To transport proteins across membranes
\end{itemize}
\begin{enumerate}
\item \textbf{Which of the following processes is most directly affected by a defect in mitochondrial DNA?}
\end{enumerate}
\begin{itemize}
\item Glycolysis
\item ATP production through oxidative phosphorylation
\item Protein synthesis in the cytoplasm
\item Lipid metabolism in the smooth ER
\item DNA replication in the nucleus
\end{itemize}
\begin{enumerate}
\item \textbf{Which of the following molecules is required for the initiation of translation in eukaryotic cells?}
\end{enumerate}
\begin{itemize}
\item DNA polymerase
\item RNA polymerase
\item Initiator tRNA (Met-tRNAi)
\item Ribosomal subunits
\item Elongation factors
\end{itemize}
\begin{enumerate}
\item \textbf{In the absence of oxygen, how does the cell regenerate NAD+ to allow glycolysis to continue?}
\end{enumerate}
\begin{itemize}
\item By oxidizing NADH in the electron transport chain
\item By reducing pyruvate to lactate through fermentation
\item By transferring electrons to FADH2
\item By generating ATP through substrate-level phosphorylation
\item By synthesizing new NAD+ molecules in the cytoplasm
\end{itemize}
\begin{enumerate}
\item \textbf{Which cellular process is most likely disrupted by defects in clathrin-mediated endocytosis?}
\end{enumerate}
\begin{itemize}
\item Exocytosis of neurotransmitters
\item Uptake of specific molecules through receptor-mediated endocytosis
\item Recycling of plasma membrane proteins
\item Mitochondrial fission and fusion
\item Formation of secretory vesicles in the Golgi apparatus
\end{itemize}
\begin{enumerate}
\item \textbf{Which of the following molecules acts as a second messenger in intracellular signaling pathways?}
\end{enumerate}
\begin{itemize}
\item ATP
\item cAMP (cyclic AMP)
\item GTP
\item DNA
\item RNA
\item \textbf{Which of the following best describes the process of autophagy in cells?}
\end{itemize}
\begin{itemize}
\item Degradation of damaged organelles within lysosomes
\item Secretion of proteins from the Golgi apparatus
\item Formation of vesicles at the plasma membrane
\item Endocytosis of large particles such as bacteria
\item Transfer of proteins between the rough ER and the Golgi apparatus
\end{itemize}
\begin{enumerate}
\item \textbf{Which enzyme is primarily responsible for converting glucose to glucose-6-phosphate in the first step of glycolysis?}
\end{enumerate}
\begin{itemize}
\item Hexokinase
\item Phosphofructokinase
\item Glucose-6-phosphatase
\item Pyruvate kinase
\item Aldolase
\end{itemize}
\begin{enumerate}
\item \textbf{Which of the following mechanisms is involved in the regulation of gene expression at the post-transcriptional level?}
\end{enumerate}
\begin{itemize}
\item Histone acetylation
\item DNA methylation
\item RNA splicing
\item Promoter binding by transcription factors
\item Chromatin remodeling
\end{itemize}
\begin{enumerate}
\item \textbf{Which of the following steps in the TCA cycle directly produces a molecule of GTP (or ATP)?}
\end{enumerate}
\begin{itemize}
\item Conversion of citrate to isocitrate
\item Conversion of succinyl-CoA to succinate
\item Oxidation of malate to oxaloacetate
\item Conversion of fumarate to malate
\item Oxidation of alpha-ketoglutarate to succinyl-CoA
\end{itemize}
\begin{enumerate}
\item \textbf{Which of the following types of cellular junctions is primarily involved in providing mechanical strength to tissues, such as the skin?}
\end{enumerate}
\begin{itemize}
\item Tight junctions
\item Desmosomes
\item Gap junctions
\item Adherens junctions
\item Hemidesmosomes
\end{itemize}
\begin{enumerate}
\item \textbf{What is the role of NAD+ in cellular respiration?}
\end{enumerate}
\begin{itemize}
\item To act as the final electron acceptor in the electron transport chain
\item To transport protons across the mitochondrial membrane
\item To carry electrons from glycolysis and the TCA cycle to the electron transport chain
\item To facilitate the synthesis of ATP in the cytoplasm
\item To catalyze the conversion of pyruvate to acetyl-CoA
\end{itemize}
\begin{enumerate}
\item \textbf{Which of the following molecules is most important for the elongation phase of protein synthesis?}
\end{enumerate}
\begin{itemize}
\item DNA polymerase
\item RNA polymerase
\item Aminoacyl-tRNA synthetase
\item Elongation factor Tu (EF-Tu)
\item Signal recognition particle (SRP)
\end{itemize}
\begin{enumerate}
\item \textbf{Which of the following processes is primarily responsible for generating proton gradients across membranes in the mitochondria?}
\end{enumerate}
\begin{itemize}
\item Glycolysis
\item The TCA cycle
\item The electron transport chain
\item Oxidative phosphorylation
\item Beta-oxidation
\end{itemize}
\begin{enumerate}
\item \textbf{Which of the following mechanisms is responsible for activating zymogens, such as trypsinogen, into their active forms?}
\end{enumerate}
\begin{itemize}
\item Proteolytic cleavage
\item Phosphorylation
\item Ubiquitination
\item Allosteric regulation
\item Methylation\\ 
\end{itemize}
CHATGPT ANSWERS\par
\begin{enumerate}
\item \textbf{Their DNA is organized into circular chromosomes within the nucleus.}
\item \textbf{By using a pinhole to eliminate out-of-focus light and capturing light from a single focal plane.}
\item \textbf{In the extracellular matrix.}
\item \textbf{Breakdown of macromolecules.}
\item \textbf{Fatty acid.}
\item \textbf{By spanning the lipid bilayer with one or more hydrophobic regions.}
\item \textbf{Processing of N-linked carbohydrates.}
\item \textbf{3.}
\item \textbf{Resilience to compression forces.}
\item \textbf{To initiate the branching of actin filaments.}
\item \textbf{Activation of a receptor leads to downstream effects on transcription factors.}
\item \textbf{Intracellular receptor.}
\item \textbf{Cells will contain multiple nuclei.}
\item \textbf{Cancer cells continue to divide even when they are tightly packed together.}
\item \textbf{Increased cell production in normal tissue.}
\item \textbf{Cell shrinkage.}
\item \textbf{Self-renewal.}
\item \textbf{Paracrine signaling.}
\item \textbf{A = proteoglycan, B = glycosaminoglycan, C = repeating disaccharide.}
\item \textbf{(1) Hydrogen bond; (2) Ionic bond; (3) Disulfide bond; (4) Peptide bond.}
\item \textbf{Apoenzyme and co-factor join together to form a functional enzyme.}
\item \textbf{Substrate concentration [S].}
\item \textbf{All options.}
\item \textbf{Riboflavin.}
\item \textbf{Fructose 1,6-bisphosphate}
\item \textbf{Oxygen}
\item \textbf{Transcription of messenger RNA (mRNA)}
\item \textbf{To degrade or prevent the translation of target mRNAs}
\item \textbf{It prevents the shortening of telomeres by adding repetitive sequences to chromosome ends}
\item \textbf{Helicase}
\item \textbf{Activation of adenylyl cyclase and production of cyclic AMP (cAMP)}
\item \textbf{Phosphorylation of retinoblastoma (Rb) protein}
\item \textbf{Frameshift mutation}
\item \textbf{To remove introns and join exons together}
\item \textbf{Euchromatin is loosely packed and associated with active gene expression, while heterochromatin is densely packed and transcriptionally silent}
\item \textbf{To catalyze the folding of proteins into their functional forms}
\item \textbf{ATP production through oxidative phosphorylation}
\item \textbf{Initiator tRNA (Met-tRNAi)}
\item \textbf{By reducing pyruvate to lactate through fermentation}
\item \textbf{Uptake of specific molecules through receptor-mediated endocytosis}
\item \textbf{cAMP (cyclic AMP)}
\item \textbf{Degradation of damaged organelles within lysosomes}
\item \textbf{Hexokinase}
\item \textbf{RNA splicing}
\item \textbf{Conversion of succinyl-CoA to succinate}
\item \textbf{Desmosomes}
\item \textbf{To carry electrons from glycolysis and the TCA cycle to the electron transport chain}
\item \textbf{Elongation factor Tu (EF-Tu)}
\item \textbf{The electron transport chain}
\item \textbf{Proteolytic cleavage}
\end{enumerate}
\end{document}